\chapter{Evaluation} \label{sec:Evaluation}

To evaluate the usability of the Praxsmith project, two methods were used:

\begin{itemize}
    \item We used the tool from the perspective of developers who were already familiar
    with the syntax, technologies, and surrounding literature.
    \item A trial was conducted with other individuals, involving a set of surveys
    and experimentation with the language and engine taking place in between. This
    survey was flawed in multiple ways, despite providing useful feedback, as detailed
    in \Cref{sec:Limitations}. It should be treated more as exploratory evidence
    from simple case studies.
\end{itemize}

\section{Research Protocol} \label{sec:ResearchProtocol}

All materials used during the recruitment and survey process can be found in
\Cref{sec:SurveyMaterials}.

First, a recruitment post was distributed to a few select places:

\begin{itemize}
    \item Friends of the authors.
    \item The Cal Poly Game Development Club Discord server.
    \item The Cal Poly Gender Inclusive Gaming Club Discord server.
\end{itemize}

Users who expressed interest were able to fill out a form containing a number
of questions intended to gauge their interest and experience in games and game
development. Part of this form was a preferred email for communication, which
was then used to send more information about the survey, plus a link to the
project itself. Participants were asked to spend as much time as they felt
comfortable developing, which ended up averaging around 1 hour.

Finally, participants were asked to complete a post-experience survey detailing
their experience with the project. This included Likert-scale questions, multiple-choice
questions, and text-entry questions. Four individuals completed this final step,
providing real-world information about how a newcomer might approach a tool like this.

All survey materials and procedures were approved by the Cal Poly San Luis Obispo IRB
(IRB Protcol \#2026-080).

\section{Engine Capabilities} \label{sec:EngineCapabilities}

Users seemed to be satisfied with the capabilities of the language and engine. Two
expressed this directly, stating that they feel like they could do much more with
it if given the time to explore. One mentioned that the modularity of the practice
system gave them many options for setting up their scenarios and dropping characters
into different situations.

There were minimal concerns within evaluations about the
core concepts explored with Praxsmith. All participants were satisfied with their
resulting product after experimentation. The challenge then becomes adjusting how
they interact with those inner systems so that they can convey their ideas more
easily.

\section{Documentation And Examples} \label{sec:DocumentationAndExamples}

As the author of any project, it is hard to evaluate the approachability of a
system from the standpoint of someone who has never touched it. Great care was taken
in building a documentation page for the language, including tips on how to use it
and three short examples showcasing different features and potentials.

All participants had positive reviews about the documentation and examples, with
2\slash 4 participants rating the usefulness of the documentation at full score (5/5) on the
Likert scale and 3\slash 4 participants giving a full score (5/5) to the provided examples.
One user pointed out in one of their responses that a question they had showed up
in one of the examples, allowing them to move forward with their project quickly.

As a contrary point, one user expressed confusion about the purpose of the
$delete\ self$ outcome found in many of the example practices, which is a clarity
issue with the documentation. This confusion proves that although the existence of the
documentation was much appreciated, there are still holes in it that should be
identified and patched by more user studies.

\section{More Types} \label{sec:MoreTypes}

As one of the users was experimenting with an inventory system within their game,
they expressed the desire for array type support within the language. This would allow
for more complex interaction, like characters holding several items at once or
remembering a series of actors for later actions.

It is worth noting that arrays are distinct datatypes and would alter the shape
of the language. It was for this reason that they were not included, and
situations such as actors holding many things can instead be represented by an
exclusive $is\_held\_by$ relation, paired with adding dummy objects into the world.
That way, objects can only be held by one individual, allowing for possession to
be represented in a more lifelike manner.

However, this complaint highlights a shortcoming of the language either way. Either
the existing solution needs to be made more obvious and easier to work with, or
arrays should be added to the system as a way to support this functionality.
Inventory systems are incredibly common in games, so more research into this area
would be appropriate.

\section{If-Statements In Effects} \label{sec:IfStatementsInEffects}

One user noted the lack of a dedicated system for comparisons within the language.
Praxsmith provides if-statements for boolean values in the form of $and$ short
circuiting and can be extended to if\slash else by adding $or$ and $not$ statements, but
there are two main problems with this:

\begin{itemize}
    \item It is clunky. Especially once you get into if\slash else territory, the long
    expressions quickly become unreadable.
    \item It only works when returning boolean values, which is okay for conditions,
    but outcomes\slash effects suffer as a result. Changing dialog slightly based on some
    state variable is one example of where this might be useful, as that currently
    requires one action within a practice for each dialog choice, which gets
    repetitive quickly. In their post-experience survey, one survey participant pointed
    out this repetition in particular as a point of inconvenience.
\end{itemize}

This is likely the first target for language design change, as it would improve quality
of life greatly and allow for much more dynamic and expressive practice outcomes.

\section{Unique Syntax} \label{sec:UniqueSyntax}

Users of the language were universally surprised by the syntax at first, but said they
were able to adjust quickly. One user pointed out practices (e.g. $practice.greet.alice.bob$)
as a core example, noting that it became intuitive after regular usage.

This syntax is carried over from the sentences outlined in the original Praxis paper. \cite{evans_ai_2015}
Since this project no longer strictly follows the exclusion logic laid out in
Praxis, sentences could be replaced with a more appropriate representation. On the other
hand, they work nicely as a singular unit and may be the cleanest option, especially
considering users seemed to adjust quickly. More research should be put into this
possibility.

\section{User Interface} \label{sec:UserInterface}

Only two users reported minor complaints about the user interface of the project.
One was to do with the visual highlighting being unclear, and another expressed that
they would like to be able to see the scripts as they are playing the game for
reference. Given the minimal nature of these complaints, it can be concluded that
the development environment did not significantly impact user experience with the language.

\section{Verbosity} \label{sec:Verbosity}

\begin{lstlisting}[float=hb, caption=A simple Do Nothing practice., label=lst:DoNothingPractice]
practice do_nothing (Actor) {
  actions: [
    {
      for: Actor
      name: "Do nothing"
      conditions: []
      outcomes: []
    }
  ]
}
\end{lstlisting}

The language ended up being fairly verbose, which made line counts skyrocket in even
simple scenarios. \Cref{lst:DoNothingPractice} describes the simplest possible practice,
one that does nothing, and it still takes up 10 lines with legible bracket practices,
plus 4 levels of indentation. Adding conditions and outcomes, some of which take up
multiple lines, plus more indentation on their own. Before more features are added,
it seems more ideal to cut down on what is already there to make it more legible and
easier to understand, and the project would benefit from more thought in this
department.

One potential solution is to implement an authoring interface, similar to that found
in the Ensemble Engine. \cite{samuel_ensemble_2015} They note that a previous Comme il
Faut developer was able to alleviate many authoring struggles by creating a dedicated
interface for building out the world. With collapsible practices and quick overviews of
different relations, a user interface for authoring is likely a significant improvement
over raw code.

\section{Writer Friendliness} \label{sec:WriterFriendliness}

Praxis is only one part of the Versu project, the other part being a play script-like
syntax named Prompter that behaves similarly to the Inform 7 project.
\cite{nelson_prompter_2014} \cite{inform_7_inform_nodate} The lack of
this final layer of abstraction is felt when using Praxsmith. Although the user study
was not able to cover this due to selection bias (see \Cref{sec:SurveyBias}), we are
not confident that those without programming experience would be able to
easily convey their ideas given the technicality of the language. As part of the
purpose of this project was to allow writers to express their ideas more easily,
this is a significant limitation.

This upper layer of abstraction is a primary opportunity for future research.
Including shortcuts for those who aren't as comfortable with traditional
programming syntax and methods is an ideal next step.

\section{Error Handling} \label{sec:ErrorHandling}

One of the philosophies we stuck to while implementing the type checker is
that the user should never see runtime errors. Everything should either be caught
during build time or fall back to predictable default cases. Especially in a simulation
engine designed to handle heavily branching cases and unpredictable situations, the
risk of players encountering situations you never even thought to test is very high.
Logic errors are bound to occur either way, but the design of the language and a strict
type checker can both prevent common mistakes.

As far as we know, the engine succeeds in this goal, and errors were caught
as the world is built for all scenarios encountered. Appropriate context was provided
as well, with line numbers for parsing errors, thanks to the errors returned by ``pest''
and a hierarchy of information that points to relevant type errors, thanks to Rust's
``anyhow'' crate and appropriate handling within the codebase
(\Cref{fig:TypeCheckerError} and \Cref{fig:ParsingError}). Even though runtime
errors are designed not to occur, the error handling within those is just as informative
so developers can target issues quickly if they arise.

\begin{figure}
    \centering
    \includegraphics[width=6in]{type_checker_error}
    \caption{An error reported to the user during type checking.}
    \label{fig:TypeCheckerError}
\end{figure}

\begin{figure}
    \centering
    \includegraphics[width=6in]{parsing_error}
    \caption{An error reported to the user during parsing.}
    \label{fig:ParsingError}
\end{figure}

This careful handling is felt throughout development, and quick feedback was invaluable
to the iteration process. Instead of realizing a guarantee was not provided for one
rare practice deep into playtesting or not at all, the engine forces you to provide
a fallback case and clarify your intentions.

\section{Actor Turns} \label{sec:ActorTurns}

One topic that we did not find commonly discussed in papers reviewed was the topic of
turns during interactions. This is easy to gloss over until implementation, when it
becomes unavoidable. In a conversation, people don't go in a circle. The conversation
flows naturally (ideally), and people speak when they feel it is appropriate. If someone
is complimented, they respond immediately instead of waiting for everyone else's turn
to speak.

Praxsmith took the oversimplified approach of a round-robin style turn-based system.
This was easiest to implement and led to predictable outcomes, but it did come with
the caveats above.

One potential solution is turn forcing. After a compliment, the turn could be manually
redirected to the person receiving the compliment, which forces them to address it
without introduction. This also opens up the door to allow certain actors, like managers,
to perform many actions at a time, which helps reduce the amount of oversight required.
The Ensemble Engine implementation behaves similarly, where a character is expected to
respond to a social exchange immediately when prompted. \cite{samuel_ensemble_2015}

Any proposed solution requires integrating turns more closely with the engine, which
we intentionally avoided because it places more authority on the one using
the engine. For example, a Unity integration would allow characters to interact in
physical space with one another, which would trigger actions; however, moving action
order to the engine level seems to be worth looking into for future research.

\section{Value Decay} \label{sec:ValueDecay}

When developing a more complex scenario, we wanted to create a situation
in which characters can perform actions to feel $proud$, and as a player you can
encourage them to take different actions to achieve this goal throughout the story.
To make this happen, there needed to be some way for the $proud$ trait to ``decay'',
so that their goal to become $proud$ wasn't already satisfied at the beginning of the
turn. An increasing number field would work, but only until it hits max, and it wouldn't
be a very intuitive solution. Another solution could be creating a ``pride manager''
actor for each character that automatically removes the character's pride every turn, but
that requires a lot of scaffolding, especially if you start extending the decay over
multiple turns.

Another example could be a $greet$ practice, which is only relevant to acknowledge
the turn after the greeting is made. It would be awkward to wait several exchanges
into a conversation before returning the greeting.

The Ensemble Engine supplies a way to assign a duration to certain facts within the
world. \cite{samuel_ensemble_2015} Some similar mechanism by which relations (and
potentially their fields) can decay would be worth looking into for a future iteration
of this project.

\section{Other Language Tweaks} \label{sec:OtherLanguageTweaks}

There are a few other inconveniences in the language that become apparent when using it
for an extended duration:

\begin{itemize}
    \item Nesting operators that introduce bindings (e.g., $count\ A,\ B\ where\ A.likes.B$)
    turns out to be a common pattern, especially when working with goals. Currently,
    this cannot be expressed in a single expression and must be broken up as in
    $sum\ (count\ B\ where\ A.likes.B)\ across\ A\ where\ @A\ is\ @A$. Support for the
    former syntax would make lines like these much more manageable.
    \item Where clauses in aggregators (e.g. $sum\ \_\ across\ \_\ where\ \_$) turn out
    to be unnecessary in some nested scenarios, as in the above example. Ternary
    if-statements would alleviate their purpose entirely and avoid cluttering
    the language with single-purpose keywords.
    \item Agents are sometimes referred to with a $@$ prefix and sometimes without, which
    can be confusing. Normalizing this behavior would be a welcome change.
    \item Being able to generate multiple actions for any actors that fit the conditions
    would be very useful, especially for manager practices. For example, a practice that
    decays any $pride$ field over time, as described in \Cref{sec:ValueDecay}.
\end{itemize}

\section{Actor Individuality} \label{sec:ActorIndividuality}

In a framework where interactions are governed by generic social practices, individuality
was a concern of ours going into development. However, during testing sessions,
Being able to provide custom dialog as fields (e.g. a custom greeting to a $greet$
practice) ended up being a clean solution to this. More research should be done into
how this feels to use at a larger scale, but it feels like an appropriate solution
that emerged from the design of the language.

\section{Debugger} \label{sec:Debugger}

One thing the Ensemble Engine developers emphasized is their debugging utilities. \cite{samuel_ensemble_2015} The Praxsmith
web interface has a simplified version of this, where developers can evaluate expressions
and run effects in real-time, plus see a list of all active relationships and their
fields. However, there is much more potential for this, especially with the graph-based
approach the engine internals already use (see \Cref{sec:InternalGraphRepresentation}). Adding more debugging capability
to the interface would make the development experience much smoother and would complement
the build-time guarantees nicely.
